% part: intuitionistic-logic
% chapter: propositions-as-types
% section: terms-to-proofs

\documentclass[../../../include/open-logic-section]{subfiles}

\begin{document}

\olfileid{int}{pty}{tp}

\olsection{Recovering \usetoken{P}{derivation} from Proof Terms}

Now let us consider the other direction: translating a proof term~$N$
back to !!a{proof} in \Log{N2i}. In general, it is only possible to do
this relative to a context that contains pairs $x : !A$ for every
variable~$x$ that occurs free in the given proof term. But let's start
with an example where there are no free variables, namely the term
\[
\lambd[\typeof{x}{!A \land
  !B}][\pair{\proj{2}{x}}{\proj{1}{x}}]
\]
It is of the form $\lambd[\typeof{x}{!A \land !B}][N_1]$. Since the
only rule of \Log{tN2i} that produces a term of this form
is~\Intro{\lif}, we know that the !!{proof} encoded by~$N$ has to end
with~\Intro{\lif}, i.e., it ends in
  \[
  \Axiom$x : !A\land !B \fCenter \pair{\proj{2}{x}}{\proj{1}{x}} : 
    !C$
  \RightLabel{\Intro\lif}
  \UnaryInf$\fCenter \lambd[\typeof{x}{!A \land
  !B}][\pair{\proj{2}{x}}{\proj{1}{x}}] : (!A \land !B) \lif !C$
  \DisplayProof
  \]
We don't yet know what $!C$ is, but we do know that the antecedent is
$!A \land !B$, and that the context of the premise has to include
$x:!A \land !B$.

Now the proof term associated with the premise is already determined:
$\pair{\proj{2}{x}}{\proj{1}{x}}$. The !!{proof} \emph{it} encodes
must end with \Intro{\land}. We still don't know what $!C$ is, but
because it is the conclusion of~\Intro{\land}, we know it must be a
conjunction, i.e., $!C \ident (!C_1 \land !C_2)$. We continue the
reconstruction:
  \[
  \Axiom$x : !A\land !B \fCenter \proj{2}{x} : !C_1$
  \Axiom$x : !A\land !B \fCenter \proj{1}{x} : !C_2$
  \RightLabel{\Intro\land}
  \BinaryInf$x : !A\land !B \fCenter \pair{\proj{2}{x}}{\proj{1}{x}} : 
    !C_1 \land !C_2$
  \RightLabel{\Intro\lif}
  \UnaryInf$\fCenter \lambd[\typeof{x}{!A \land
  !B}][\pair{\proj{2}{x}}{\proj{1}{x}}] : (!A \land !B) \lif (!C_1 \land !C_2)$
  \DisplayProof
  \]
The proof term $\pi_2(x)$ means that top sequent on the left was
obtained by \Elim{\land}[2], i.e., the premise is of the form $!D
\land !C_1$. On the right, we can determine that the inference leading
to $!C_2$ must be \Elim{\land}[1], and that the premise is of the form
$!C_2 \land !E$.
  \[
  \Axiom$x : !A\land !B \fCenter x: !D\land !C_1$
  \RightLabel{\Elim\land[2]}
  \UnaryInf$x : !A\land !B \fCenter \proj{2}{x} : !C_1$
  \Axiom$x : !A\land !B \fCenter x: !C_2\land !E$
  \RightLabel{\Elim\land[1]}
  \UnaryInf$x : !A\land !B \fCenter \proj{1}{x} : !C_2$
  \RightLabel{\Intro\land}
  \BinaryInf$x : !A\land !B \fCenter \pair{\proj{2}{x}}{\proj{1}{x}} : 
    !C_1 \land !C_2$
  \RightLabel{\Intro\lif}
  \UnaryInf$\fCenter \lambd[\typeof{x}{!A \land
  !B}][\pair{\proj{2}{x}}{\proj{1}{x}}] : (!A \land !B) \lif (!C_1 \land !C_2)$
  \DisplayProof
  \]
Since the proof terms of the topmost sequents are now variables, we
know that they must be axioms. This determines $!C_1$, $!C_2$, $!D$, and
$!E$: $!D \ident !A$ and $!C_1 \ident !B$ or else the left sequent is
not an axiom, and $!C_2 \ident !A$ and $!E \ident !B$, or else the
right sequent isn't. We have recovered the proof fully:
  \[
  \Axiom$x : !A\land !B \fCenter x: !A\land !B$
  \RightLabel{\Elim\land[2]}
  \UnaryInf$x : !A\land !B \fCenter \proj{2}{x} : !B$
  \Axiom$x : !A\land !B \fCenter x: !A\land !B$
  \RightLabel{\Elim\land[1]}
  \UnaryInf$x : !A\land !B \fCenter \proj{1}{x} : !A$
  \RightLabel{\Intro\land}
  \BinaryInf$x : !A\land !B \fCenter \pair{\proj{2}{x}}{\proj{1}{x}} : 
    !B \land !A$
  \RightLabel{\Intro\lif}
  \UnaryInf$\fCenter \lambd[\typeof{x}{!A \land
  !B}][\pair{\proj{2}{x}}{\proj{1}{x}}] : (!A \land !B) \lif (!B \land !A)$
  \DisplayProof
  \]


\end{document}
